% Generated draft: requires mathematical and source review.
\documentclass[11pt]{article}
\usepackage[margin=25mm]{geometry}
\usepackage{fontspec}
\setmainfont{Latin Modern Roman}
\usepackage{amsmath,amssymb,mathtools,unicode-math}
\setmathfont{Latin Modern Math}
\usepackage{enumitem,needspace,xurl,hyperref}
\setlength{\parindent}{0pt}
\setlength{\parskip}{5pt}
\setlength{\emergencystretch}{3em}
\newcommand{\N}{\mathbb{N}}
\newcommand{\Z}{\mathbb{Z}}
\newcommand{\Q}{\mathbb{Q}}
\newcommand{\R}{\mathbb{R}}
\newcommand{\C}{\mathbb{C}}
\newcommand{\F}{\mathbb{F}}
\newcommand{\A}{\mathbb{A}}
\newcommand{\PP}{\mathbb{P}}
\newcommand{\E}{\mathbb{E}}
\newcommand{\eps}{\varepsilon}
\newcommand{\dd}{\,\mathrm{d}}
\begin{document}
\section*{UnsolvedMath ID: 10400190}
\textbf{Source problem code:} catalogue entry
\subsection{Definitions and mathematical statement}
Let \(\mathcal H\) be the set of integral homology three-spheres, with the orientation convention of the source. Let \(\sim_{Y_k}\) be its \(Y_k\)-equivalence relation, and put
\[G_{2d}=\{M\in\mathcal H:M\sim_{Y_{2d}}S^3\}/\!\sim_{Y_{2d+1}}.\]
The quotient is intended as the corresponding group under connected sum, with identity the class of \(S^3\). The conjecture states that \(G_{2d}\) is torsion free for each allowed \(d\):
\[\forall x\in G_{2d}\ \forall n\in\Z_{\geq1},\quad nx=0\Longrightarrow x=0.\]
Equivalently, for each such \(d\), each \(M\sim_{Y_{2d}}S^3\) and each positive integer \(n\),
\[\underbrace{M\mathbin{\#}\cdots\mathbin{\#}M}_{n\text{ copies}}
\sim_{Y_{2d+1}}S^3
\quad\Longrightarrow\quad M\sim_{Y_{2d+1}}S^3.\]
Here an integral homology three-sphere is a closed connected three-manifold with integral homology isomorphic to that of \(S^3\). The research summary identifies this as the intended domain, which is not written in the short formula itself.

In the customary clasper setting, \(Y_k\)-equivalence is generated by surgeries on suitable degree-\(k\) tree claspers (with filtration and degree conventions fixed in the source). This excerpt supplies no definition, degree convention or proof of the group structure, so these must be supplied from the preceding discussion. The range of \(d\), including whether zero is intended, is unstated. The assertion concerns the displayed filtration piece, not the entire quotient of all homology spheres by \(Y_{2d+1}\)-equivalence.
\subsection{Short English statement}
No nonzero class in the indicated even-degree \(Y\)-filtration quotient of homology spheres should have finite additive order. Equivalently, a finite connected sum becoming trivial in the next quotient should force the original class to be trivial there.
\subsection{Sources}
\begin{itemize}
\item Local UnsolvedMath registry, permanent ID \(10400190\), exact source code \texttt{catalogue entry}.
\item Ohtsuki, \emph{Problems on invariants of knots and 3-manifolds} (2002), Conjecture 10.15, PDF page 152: \url{https://msp.org/gtm/2002/04/gtm-2002-04-024s.pdf}. Recorded access 2026-07-29; extraction pdf-font-structure.
\end{itemize}
\end{document}
